\par\Needspace{18\baselineskip}\bigskip
\pdfbookmark[2]{Section première. Mouvement elliptique.}{pallas-II-section1}
\begin{center}
{\large\bfseries Section première.\par}
\medskip\textit{Mouvement elliptique.}
\end{center}
\begin{flushright}
\begin{minipage}{.94\textwidth}
\small\raggedright
. . . . . Ye other wandring fires that move\\
In mystic dance, not without song, resound\\
His praise, who out of darkness call’d up light.\\
\hspace*{2em}\textsc{Milton}, Paradise lost.
\end{minipage}
\end{flushright}
\pdfbookmark[3]{Exposition, 1}{pallas-II-article1}
\begin{center}1.\end{center}
\noindent\hspace*{\parindent}Soient \(x,y,z\) les coordonnées d’une planète relativement à trois plans perpendiculaires entre eux passant par le centre du soleil; soit de plus \(t\) le tems, \(m\) la somme des masses du soleil et de la planète, \(r\) la distance de la \SourceLeaf{447}planète au soleil. Ainsi en faisant abstraction des perturbations nous aurons
\[
\begin{array}{rl}
1)\quad&0=\dfrac{\mathrm{dd}x}{\mathrm{d}t^2}+\dfrac{mx}{r^3}\\[5pt]
2)\quad&0=\dfrac{\mathrm{dd}y}{\mathrm{d}t^2}+\dfrac{my}{r^3}\\[5pt]
3)\quad&0=\dfrac{\mathrm{dd}z}{\mathrm{d}t^2}+\dfrac{mz}{r^3}\mathpunct{\cdot}
\end{array}
\]
\par Il s’ensuit, que faisant
\[
\begin{array}{rl}
4)\quad&\dfrac{y\mathrm{d}z-z\mathrm{d}y}{\mathrm{d}t}=A\\[5pt]
5)\quad&\dfrac{z\mathrm{d}x-x\mathrm{d}z}{\mathrm{d}t}=B\\[5pt]
6)\quad&\dfrac{x\mathrm{d}y-y\mathrm{d}x}{\mathrm{d}t}=C,
\end{array}
\]
\(A,B,C\) seront des quantités constantes, qui par conséquent dépendront de l’état primitif. On tire de ces équations
\[
\begin{array}{rl}7)\quad&Ax+By+Cz=0.\end{array}
\]
\par L’orbite sera donc dans un plan passant par le centre du soleil. Nous avons de plus
\[
\begin{gathered}
rr=xx+yy+zz\\
r\mathrm{d}r=x\mathrm{d}x+y\mathrm{d}y+z\mathrm{d}z\\
\begin{aligned}
r\mathrm{dd}r+\mathrm{d}r^2&=x\mathrm{dd}x+y\mathrm{dd}y+z\mathrm{dd}z+\mathrm{d}x^2+\mathrm{d}y^2+\mathrm{d}z^2\\
&=\mathrm{d}x^2+\mathrm{d}y^2+\mathrm{d}z^2-\frac{m\mathrm{d}t^2}{r}\mathpunct{\cdot}
\end{aligned}
\end{gathered}
\]
\par Or les équations 4, 5, 6 donnent
\[
\begin{aligned}
(AA+BB+CC)\mathrm{d}t^2
&=(xx+yy+zz)(\mathrm{d}x^2+\mathrm{d}y^2+\mathrm{d}z^2)\\
&\qquad{}-(x\mathrm{d}y+y\mathrm{d}y+z\mathrm{d}z)^2\\
&=rr(\mathrm{d}x^2+\mathrm{d}y^2+\mathrm{d}z^2)-rr\mathrm{d}r^2\\
&=r^3\mathrm{dd}r+mr\mathrm{d}t^2.
\end{aligned}
\]
\par Ainsi en faisant
\[
\frac{AA+BB+CC}{m}=p,\qquad r-p=u,
\]
on aura \(0=\frac{\mathrm{dd}u}{\mathrm{d}t^2}+\frac{u}{r^3}\mathpunct{\raisebox{.45ex}{,}}\) équation analogue aux trois premières. On voit donc qu’en faisant
\[
\begin{array}{rl}
8)\quad&\dfrac{x\mathrm{d}u-u\mathrm{d}x}{\mathrm{d}t}=F\\[5pt]
9)\quad&\dfrac{y\mathrm{d}u-u\mathrm{d}y}{\mathrm{d}t}=G\\[5pt]
10)\quad&\dfrac{z\mathrm{d}u-u\mathrm{d}z}{\mathrm{d}t}=H,
\end{array}
\]
\SourceLeaf{448}\(F,G,H\) seront aussi des quantités constantes, dépendantes de l’état primitif. Combinant ces formules avec 4, 5, 6 on en tire
\[
\begin{array}{rl}
11)\quad&Gz-Hy=A(r-p)\\
12)\quad&Hx-Fz=B(r-p)\\
13)\quad&Fy-Gx=C(r-p).
\end{array}
\]
\par On conclut facilement de ces équations, que les projections de l’orbite et l’orbite elle-même sont des courbes du second ordre, que l’orbite a pour demi-paramètre \(p\), et que le centre du soleil occupe un de ses foyers. Nous remarquerons encore, qu’elles donnent tout de suite une équation de condition entre les six constantes \(A,B,C,F,G,H\)
\[
\begin{array}{rl}14)\quad&AF+BG+CH=0.\end{array}
\]
\par De même on tire des équations 8, 9, 10
\[
\begin{aligned}
&(FF+GG+HH)\mathrm{d}t^2\\
&=(xx+yy+zz)\mathrm{d}u^2-2(x\mathrm{d}x+y\mathrm{d}y+z\mathrm{d}z)u\mathrm{d}u\\
&\qquad{}+(\mathrm{d}x^2+\mathrm{d}y^2+\mathrm{d}z^2)uu\\
&=rr\mathrm{d}r^2-2r\mathrm{d}r\mathpunct{.}(r-p)\mathrm{d}r+(\mathrm{d}x^2+\mathrm{d}y^2+\mathrm{d}z^2)(r-p)^2\\
&=(AA+BB+CC)\left(1-\frac{2p}{r}\right)\mathrm{d}t^2+(\mathrm{d}x^2+\mathrm{d}y^2+\mathrm{d}z^2)pp.
\end{aligned}
\]
\par Donc en faisant
\[
\begin{array}{rl}15)\quad&\dfrac{FF+GG+HH}{AA+BB+CC}=1-\dfrac{p}{a}\mathpunct{\raisebox{.45ex}{,}}\end{array}
\]
on aura
\[
\begin{array}{rl}16)\quad&0=\dfrac{\mathrm{d}x^2+\mathrm{d}y^2+\mathrm{d}z^2}{\mathrm{d}t^2}-\dfrac{2m}{r}+\dfrac{m}{a}\mathpunct{\cdot}\end{array}
\]
\par On aurait pu trouver cette même équation, en multipliant les équations 1, 2, 3 par \(2\mathrm{d}x,\ 2\mathrm{d}y,\ 2\mathrm{d}z\) et intégrant l’aggrégat: le développement précédent a l’avantage de montrer la liaison entre la constante, que cette intégration aurait introduite et celles qu’on avait déjà trouvées.
\par
